\scanpage{115}
% archive foliation: 114

\uline{Cor.} 5.4. Sur un corps $k$, si $X$ est lisse, séparé, de t.f., et si $F$ est \deleted{\unsure{un}} str$^{t}$ spécial de poids $j$ (resp. eff$^{t}$ adm. de poids $j$, si l'on admet conjecture de $\frac12$ simplicité), alors $R^if_*(F)$ est de eff$^{t}$ admissible de poids entre $i+j$ et $2n+2j$, et les niveaux correspondants étant resp. $\leq\unsure{2(n-i+j)|2n-i+j)},\unsure{2n-i+j-1|2(n-i+j-1},\ldots,0$. Si on a $i\leq n$, on a \deleted{les niveaux}\inserted{seulement} des poids \deleted{\unsure{et la suite de}} $i+j,i+j+1,\ldots,2(i+j)$, \deleted{\unsure{de}} niveaux correspondants étant $\leq i+j,i+j-1,\ldots,0$.

\margin{résolution, Lefschetz fort, et le cas échéant semi-simplicité}

On peut, soit procéder par dualité à partir de 5.1 et 5.2, soit reprendre \inserted{dans le cas $F$ quelc.} la dém. de 5.1 resp. 5.2, en utilisant la \deleted{suite exacte de cohomologie relative} \inserted{\unsure{trivialité|transitivité} de la suite spectrale de Leray}, pour se ramener au cas où $F$ est $\mathbb Q_l$. Dans ce cas, on peut procéder par \inserted{compactification et} résolution des singularités, pour écrire $X=\widehat X-D$, $\widehat X$ \deleted{projectif et lisse} propre et lisse et $D$ \inserted{div} à croisements normaux, et en écrivant suite exacte de ch. relative (supports quelc.) pour $X,\widehat X$ et $D$, et la suite spectrale de localisation pour $H^{\unsure{\alpha|*}}_D(\widehat X)$.

\margin{le faire à titre de vérification de cohérence du yoga!}

6. Soit $f:X\longrightarrow S$ un morphisme de t.f., et $F$ un faisceau \deleted{très spécial}\inserted{effectivement admissible} sur $X$, de poids $j$. On veut des renseignements sur \inserted{les} $R^\nu f_*(F)$, en particulier on veut qu'ils soient effectivement admissibles, de poids que nous préciserons. \deleted{Comme} \inserted{Supposons d'abord $F=\mathbb Q_{lX}$} pour le th. de finitude, on regarde d'abord le cas d'une immersion
